% !PRINTED 415
\weberpage{415}
in which the product sign $\Pi$ extends over all divisors $\mu$ of $n$.

From this we obtain, incidentally, a proof of the number-theoretic
theorem
\weq{4}{n = \overset{\mu}{\Sigma}\varphi(\mu),}
if we set the degrees of the functions on the right and the
left side of (3) equal to one another.

On the other hand we conclude, by \gesperrt{Gauss}'s theorem
(\S.\,2), that the coefficients of all the $X_\mu$ are \gesperrt{integers}.
For if fractional numbers occurred among them, the
product could not consist solely of integral coefficients, as by (3) and (1) it nevertheless must.

All the functions $f_n(x)$ have the divisor $x-1$, and if
we carry out the division, we obtain
\weq{5}{\frac{x^n-1}{x-1} = x^{n-1} + x^{n-2} + \cdots + x + 1.}
From this we conclude that, if $r$ is any primitive
or non-primitive $n$th root of unity, with the sole exception of $r=1$, always
\weq{6}{1 + r + r^2 + \cdots + r^{n-1} = 0,}
while for $r=1$ the sum on the left side of (6)
obviously has the value $n$.

If $n$ is a prime, then apart from 1 there are no imprimitive $n$th roots of unity, and therefore, when $n$ is a prime, which we shall indicate by writing $p$ for it,
we have
\weq{7}{X_p = x^{p-1} + x^{p-2} + \cdots + x + 1.}
Just as simply, $X_n$ can be formed when $n$ is a \gesperrt{prime power}, that is,
\[
n = p^{\pi},
\]
or, if for brevity we set $p^{\pi-1} = p'$:
\[
n = p'p.
\]
The $n$th roots of unity consist in this case of the
primitive $n$th roots of unity and of the $p'$th roots of unity,
and thus:
\weq{8}{X_n = \frac{x^{pp'}-1}{x^{p'}-1} = x^{p'(p-1)} + x^{p'(p-2)} + \cdots + x^{p'} + 1}
